\documentclass[10pt]{article}
\usepackage[margin=0.6in]{geometry}
\usepackage{titlesec}
\usepackage{enumitem}
\usepackage{amsmath}
\usepackage{amssymb}
\usepackage{hyperref}
\usepackage{booktabs}
\usepackage{parskip}
\usepackage{float}
\usepackage[T1]{fontenc}

\titlespacing*{\section}{0pt}{6pt}{3pt}
\setlist[itemize]{leftmargin=1.4em, itemsep=0pt, topsep=1pt, parsep=0pt}
\setlength{\parskip}{3pt}

\title{\vspace{-2em}\large Assignment 1}
\author{\normalsize Priyank \\ \footnotesize GV~271: Graphics and Visualization, Indian Institute of Science}
\date{\footnotesize \today}

\begin{document}
\begin{center}
{\large\textbf{Assignment 1}}\\[2pt]
{\normalsize Priyank}\\
{\footnotesize GV~271 - Graphics and Visualization, Indian Institute of Science \quad\textbullet\quad \today}
\end{center}
\vspace{2pt}

\textbf{GPU:} Mesa Intel(R) Iris(R) Xe Graphics (TGL GT2)\\
\textbf{CPU:} 11th Gen Intel(R) Core(TM) i7-1165G7 @ 2.80GHz

\small
This assignment renders a protein molecule whose structure is defined by the provided OFF files. 
The mesh rotates continuously  about a randomly chosen axis with a fixed angular speed.


\section{Task 1: Model Loading, Camera Setup, Rotation, and Base Rendering}
\textbf{(i) Approach.} The OFF parser reads vertices/polygon indices into a flat indexed triangle buffer. 
Since OFF carries no normals, per vertex normals are derived from the face normals of the incident 
triangles, under either of two weighting schemes selectable at runtime (area or angle; see (iii)). 
The mesh rotates continuously about an axis drawn uniformly at random once per 
run (\texttt{initRandomAxis}), 
advancing every frame at a fixed rate driven by wall clock time 
(\texttt{timer}), composed into the model matrix with the centering 
translation and normalizing scale.

\textbf{(ii) Core files/functions.} \texttt{main.cpp}: \texttt{readOffFile}, 
\texttt{createVertexBuffers}, \texttt{createVertexNormals}, 
\texttt{ComputeModelTransform}, \texttt{setupCamera}, 
\texttt{updateProjection}, \texttt{initRandomAxis}, 
\texttt{timer}. \texttt{shader.vs}/\texttt{shader.fs} 
(\texttt{shaderProgram}) render the whole mesh.

\textbf{(iii) Challenges / design decisions.} Two vertex normal weighting
schemes are implemented and switchable at runtime (\texttt{n}, via
\texttt{currentNormalMode}/\texttt{rebuildNormals}). Accumulating
\emph{unnormalized} cross products gives area weighting for free, since a
cross product's magnitude is twice the triangle area; the default instead
normalizes each face normal and weights it by the interior angle at the
vertex, which is insensitive to triangulation density, so a dense patch of
small triangles cannot outvote a coarser neighbouring region sharing the
same vertex. Face normals are retained separately (\texttt{faceNormals})
and drive a flat shading mode for comparison. A random rotation axis can, by
chance, land close to the camera's fixed view direction, under which rotation
produces almost no visible depth change and the object appears to spin flat
rather than tumble in 3D; candidate axes within $32^\circ$ of the view direction
are rejected and redrawn.

\section{Task 2: Depth Based Colormap Shading}
\textbf{(i) Approach.} Lighting uses 
(ambient + diffuse + Phong(R.V)). The 
diffuse base color comes from a depth driven colormap 
(viridis or coolwarm) instead of a flat material, so 
once the mesh is cut open in Tasks 3 and 4, interior faces 
are distinguished by depth without needing a second material.

\textbf{(ii) Core files/functions.} \texttt{shader.fs}: 
\texttt{colormapViridis}, 
\texttt{colormapCoolWarm}, 
lighting in \texttt{main()}. 
\texttt{shader.vs}: 
\texttt{vDepth01} from \texttt{-viewPos.z}. 
Uniforms \texttt{uKa/uKd/uKs/uShininess}, 
\texttt{uNearDepth/uFarDepth} set once, 
Functions : 
\texttt{faceNormals},
\texttt{rebuildNormals()},
\texttt{buildFlatShadedMesh()},
in \texttt{main.cpp}.

\textbf{(iii) Challenges / design decisions.} Lighting was evaluated 
per fragment in shaders rather than on the CPU specifically for performance: 
every quantity it depends on (rotated normal, view vector, depth color) 
changes every frame under continuous rotation regardless of implementation, 
and the GPU evaluates this in parallel across thousands of cores, 
avoiding both a serial per pixel CPU loop and again uploading a full color buffer every frame. 
Depth uses view-space
\texttt{-viewPos.z} (not NDC $z$) so the ramp stays linear regardless of projection
type. Colormaps are piecewise-linear keyframe tables, trading exact colour fidelity
for shader simplicity. A \texttt{uUseVertexColor} toggle isolates shading bugs from
geometry bugs during debugging.


\section{Task 3: CPU Side Slicing}
\textbf{(i) Approach.} $M$ evenly spaced lines through the origin ($\theta_i=i\pi/M$) 
give $2M$ rays around the full circle; each wedge is 
bounded by two half planes from its boundary rays 
rotated $\pm90^\circ$. Every triangle is clipped 
against both half planes in sequence (a Sutherland Hodgman 
clip specialized to triangles), and surviving geometry 
per wedge is pushed outward along its bisector 
for the exploded view; all buffers rebuild every 
frame. At $M=1$ the two wedges are further distinguished 
by polygon mode: one \texttt{GL\_FILL}, 
one \texttt{GL\_LINE}: satisfying the base solid/wireframe 
requirement; at $M>1$ all wedges are filled.


\textbf{(ii) Core files/functions.} \texttt{main.cpp}: \texttt{rebuildWedgeGeometryDefs}
(angles, half plane normals, bisectors), \texttt{clipListAgainstPlane} (lone vertex
detection, edge intersection, 1 or 2 triangle output), \texttt{lerpVert},
\texttt{buildAllWedgesCPU}, \texttt{uploadSliceVAO}. The $M=1$ solid/wireframe pair
is applied in \texttt{display} (the \texttt{wireframeThis} branch around
\texttt{glPolygonMode}).

\textbf{(iii) Challenges / design decisions.} The ``lone'' vertex is found without 
sorting: with three booleans and two possible values, at least two must agree, 
so a pigeonhole case split locates the odd one out in $O(1)$. The 
intersection parameter $t=d_a/(d_a-d_b)$ is guaranteed in $(0,1)$ once the 
lone vertex's sign is confirmed opposite its neighbors. Indices are cyclically relabeled 
($a,b,c=\text{lone},(\text{lone}{+}1)\%3,(\text{lone}{+}2)\%3$), preserving 
winding so lighting/culling stay correct post clip. The cost is performance: 
the full mesh is reclipped and every wedge's VBO/EBO reuploaded every frame, 
motivating Task 4.

The clipper is parameterized by signed distance $d = \mathbf{n}\cdot\mathbf{p} - w$
rather than assuming planes through the origin, so plane position is a runtime
scalar. Setting $w \neq 0$ (\texttt{o}, via \texttt{wedgeCutOffset}) slides each
wedge boundary off the scene centre in both the CPU and GPU paths. Note that 
wedges presuppose concurrency: with the boundaries moved apart,
adjacent wedges no longer share a cut line, so small gaps open between them: a
property of the wedge decomposition, not of the clip.


\section{Task 4: GPU-Side Slicing via Geometry Shader}
\textbf{(i) Approach.} The same clip algorithm runs as a 
per triangle geometry shader stage in local space, 
reusing one static buffer every frame with no CPU rebuild. 
A pass through vertex shader forwards local position/normal; 
the geometry shader clips against \texttt{uWedgeNormalA}, 
then \texttt{uWedgeNormalB} (two chained half space clips = the wedge), 
applying MVP and the offset only to emitted vertices. 
The Task 3 fill/wireframe rule at $M=1$ is 
mirrored here via \texttt{glPolygonMode} around the GPU draws.


\textbf{(ii) Core files/functions.} \texttt{slice.vs} (pass through). 
\texttt{slice.gs}: \texttt{clipTri}, \texttt{main()} (two stage clip chaining), 
\texttt{emitWorld}/\texttt{emitTri}. \texttt{shader.fs} reused unchanged. 
\texttt{main.cpp}: \texttt{buildPreppedWholeMesh}, uniform wiring for \texttt{sliceGSProgram}.


\textbf{(iii) Challenges / design decisions.} 
% \texttt{max\_vertices=12} covers the
% worst case: the first clip yields at most 2 triangles, each of which the second
% clip splits into at most 2 more, giving 4 triangles of 3 vertices. Each is emitted
% as an independent primitive (\texttt{EndPrimitive()} after every third vertex),
% so the declared \texttt{triangle\_strip} output carries no shared vertices --
% strips would not help here, since consecutive clipped triangles rarely share an edge.
% Normals are re-normalized after interpolation, since the linear blend of two unit
% vectors is not itself unit length.
\textbf{(iii) Challenges / design decisions.} \texttt{max\_vertices=12} covers the
worst case: the first clip yields at most 2 triangles, each of which the second clip
splits into at most 2 more, giving 4 triangles of 3 vertices. Each is emitted as an
independent primitive (\texttt{EndPrimitive()} after every third vertex), so the
declared \texttt{triangle\_strip} output carries no shared vertices. Normals are
renormalized after interpolation, since the linear blend of two unit vectors is not
itself unit length.

The key bug was variable shadowing: locally redeclaring
\texttt{uSWedgeNormalA/BLoc} during init left the globals at their
zero initialized value. Since $0$ is itself a valid uniform location, the wedge
normals were silently written into whichever uniform occupied slot~0 rather than
failing loudly, and the clip planes stayed at $(0,0,0)$ - for which
$\mathrm{dot}(p,n)>0$ is false everywhere, discarding every triangle. Relatedly,
uniform storage is per-program, not per name: lighting uniforms set only on
\texttt{shaderProgram} left the GPU sliced mesh pure black until set again on
\texttt{sliceGSProgram}.

When slices overlap, the extra cost lands on the fragment stage rather than the
geometry stage. If $d$ is small, neighbouring wedges still cover some of the same
pixels, and those pixels get shaded more than once: work that is thrown away
when a later fragment overwrites an earlier one. Both implementations pay the
same absolute cost here, since both rasterize the same overlapping geometry
through the same fragment shader. What differs is how much it matters. A CPU
frame is dominated by \texttt{buildAllWedgesCPU} reclipping and reuploading the
mesh - over three seconds per frame on \texttt{2OAR\_1}- so a sub millisecond
overdraw penalty disappears into the noise. GPU frames run in a few
milliseconds, so the same penalty is a visible fraction of the total.


\section{Performance Analysis: CPU vs.\ GPU Wedge Slicing}
\textbf{(i) Approach.} FPS is measured by accumulating frames over fixed wall-clock windows 
(\texttt{glutGet(GLUT\_ELAPSED\_TIME)} in \texttt{timer}), shown in the HUD. 
Two axes are varied: active planes $M$ (1 to 10, via \texttt{j}/\texttt{k}, calling
\texttt{rebuildWedgeGeometryDefs}; sampled at $M \in \{1,4,10\}$ for the table) and
mesh resolution (separate OFF files). Each (resolution, $M$) 
pair is logged for both CPU and GPU modes, with non parallel plane orientations.

\textbf{(ii) Core files/functions.} \texttt{main.cpp}: \texttt{timer}, \texttt{keyboard} (\texttt{j}/\texttt{k}), \texttt{display} (mode dispatch).

\textbf{(iii) Expected behavior, reasoning, and results.} CPU cost scales with 
(triangle count $\times$ $2M$ wedges), since \texttt{buildAllWedgesCPU} 
again clips and again uploads every wedge every frame. GPU cost grows far more
 gently with $M$: each wedge costs one more draw call plus inherently 
 parallel \texttt{slice.gs} work, with no reclip- \texttt{buildPreppedWholeMesh}'s buffer uploads once, 
 regardless of $M$. Higher resolution scales the CPU loop directly 
 and serially, while the GPU's per primitive invocation is  
 parallel and degrades more slowly until cores saturate. Wedge boundaries 
 crossing the mesh's dense interior force more triangles through the expensive 
 split and interpolate branch on both implementations- cheap on the GPU 
 (a few extra \texttt{EmitVertex} calls in GPU memory), costlier on the CPU 
 (more interpolation and a bigger reupload).


% \emph{Table below is a placeholder -- replace with FPS measured on the target machine.}

% \vspace{-2pt}
% \begin{table}[h!]
% \centering
% \footnotesize
% \renewcommand{\arraystretch}{0.92}
% \begin{tabular}{@{}llrr@{}}
% \toprule
% Mesh Resolution & Planes ($M$) & CPU FPS & GPU FPS \\
% \midrule
% Low  & 1  & TODO & TODO \\
% Low  & 5  & TODO & TODO \\
% Low  & 10 & TODO & TODO \\
% Med  & 1  & TODO & TODO \\
% Med  & 5  & TODO & TODO \\
% Med  & 10 & TODO & TODO \\
% High & 1  & TODO & TODO \\
% High & 5  & TODO & TODO \\
% High & 10 & TODO & TODO \\
% \bottomrule
% \end{tabular}
% \caption{FPS by mesh resolution and active cutting planes $M$ ($2M$ wedges), CPU vs.\ GPU slicing.}
% \end{table}

% Auto-generated by the benchmark sweep -- % Auto-generated by the benchmark sweep -- % Auto-generated by the benchmark sweep -- \input{benchmark_table.tex} directly.
\begin{table}[H]
\centering
\footnotesize
\renewcommand{\arraystretch}{0.92}
\begin{tabular}{@{}lrrrrr@{}}
\toprule
Mesh & Vertices & Triangles & Planes ($M$) & CPU FPS & GPU FPS \\
\midrule
\texttt{1GRM\_1} & 83673 & 167346 & 1 & 5.5 & 283.9 \\
\texttt{1GRM\_1} & 83673 & 167346 & 4 & 3.8 & 152.4 \\
\texttt{1GRM\_1} & 83673 & 167346 & 10 & 2.4 & 68.4 \\
\texttt{1GRM\_2} & 41836 & 83672 & 1 & 11.3 & 509.3 \\
\texttt{1GRM\_2} & 41836 & 83672 & 4 & 7.6 & 255.3 \\
\texttt{1GRM\_2} & 41836 & 83672 & 10 & 4.7 & 121.6 \\
\texttt{1GRM\_4} & 20918 & 41836 & 1 & 21.9 & 776.0 \\
\texttt{1GRM\_4} & 20918 & 41836 & 4 & 14.8 & 431.4 \\
\texttt{1GRM\_4} & 20918 & 41836 & 10 & 9.0 & 228.6 \\
\texttt{1GRM\_8} & 10459 & 20918 & 1 & 41.6 & 1000.0 \\
\texttt{1GRM\_8} & 10459 & 20918 & 4 & 29.0 & 715.7 \\
\texttt{1GRM\_8} & 10459 & 20918 & 10 & 16.5 & 358.1 \\
\texttt{1GRM\_16} & 5229 & 10458 & 1 & 82.9 & 1000.0 \\
\texttt{1GRM\_16} & 5229 & 10458 & 4 & 52.6 & 978.6 \\
\texttt{1GRM\_16} & 5229 & 10458 & 10 & 32.7 & 516.4 \\
\texttt{2OAR\_1} & 519221 & 1038654 & 1 & 0.8 & 74.8 \\
\texttt{2OAR\_1} & 519221 & 1038654 & 4 & 0.6 & 35.2 \\
\texttt{2OAR\_1} & 519221 & 1038654 & 10 & 0.4 & 15.0 \\
\texttt{2OAR\_2} & 259557 & 519326 & 1 & 1.7 & 125.5 \\
\texttt{2OAR\_2} & 259557 & 519326 & 4 & 1.2 & 63.7 \\
\texttt{2OAR\_2} & 259557 & 519326 & 10 & 0.7 & 28.1 \\
\texttt{2OAR\_4} & 129725 & 259662 & 1 & 3.5 & 209.7 \\
\texttt{2OAR\_4} & 129725 & 259662 & 4 & 2.4 & 120.2 \\
\texttt{2OAR\_4} & 129725 & 259662 & 10 & 1.4 & 54.7 \\
\texttt{2OAR\_8} & 64810 & 129830 & 1 & 6.9 & 362.3 \\
\texttt{2OAR\_8} & 64810 & 129830 & 4 & 4.8 & 205.7 \\
\texttt{2OAR\_8} & 64810 & 129830 & 10 & 2.9 & 97.7 \\
\texttt{2OAR\_16} & 32361 & 64914 & 1 & 13.8 & 600.0 \\
\texttt{2OAR\_16} & 32361 & 64914 & 4 & 9.5 & 364.3 \\
\texttt{2OAR\_16} & 32361 & 64914 & 10 & 5.7 & 161.2 \\
\texttt{3SY7\_1} & 221469 & 442946 & 1 & 2.0 & 141.4 \\
\texttt{3SY7\_1} & 221469 & 442946 & 4 & 1.4 & 75.0 \\
\texttt{3SY7\_1} & 221469 & 442946 & 10 & 0.9 & 32.9 \\
\texttt{3SY7\_2} & 110733 & 221472 & 1 & 4.1 & 253.9 \\
\texttt{3SY7\_2} & 110733 & 221472 & 4 & 2.9 & 131.9 \\
\texttt{3SY7\_2} & 110733 & 221472 & 10 & 1.8 & 60.7 \\
\texttt{3SY7\_4} & 55365 & 110736 & 1 & 6.6 & 431.4 \\
\texttt{3SY7\_4} & 55365 & 110736 & 4 & 5.7 & 236.5 \\
\texttt{3SY7\_4} & 55365 & 110736 & 10 & 3.5 & 107.0 \\
\texttt{3SY7\_8} & 27680 & 55368 & 1 & 16.3 & 688.6 \\
\texttt{3SY7\_8} & 27680 & 55368 & 4 & 11.0 & 402.3 \\
\texttt{3SY7\_8} & 27680 & 55368 & 10 & 6.8 & 179.5 \\
\texttt{3SY7\_16} & 13840 & 27684 & 1 & 31.7 & 958.6 \\
\texttt{3SY7\_16} & 13840 & 27684 & 4 & 21.5 & 610.0 \\
\texttt{3SY7\_16} & 13840 & 27684 & 10 & 13.1 & 298.1 \\
\bottomrule
\end{tabular}
\caption{Measured FPS across mesh resolutions and active cutting planes, CPU vs.\ GPU wedge slicing.}
\end{table}
 directly.
\begin{table}[H]
\centering
\footnotesize
\renewcommand{\arraystretch}{0.92}
\begin{tabular}{@{}lrrrrr@{}}
\toprule
Mesh & Vertices & Triangles & Planes ($M$) & CPU FPS & GPU FPS \\
\midrule
\texttt{1GRM\_1} & 83673 & 167346 & 1 & 5.5 & 283.9 \\
\texttt{1GRM\_1} & 83673 & 167346 & 4 & 3.8 & 152.4 \\
\texttt{1GRM\_1} & 83673 & 167346 & 10 & 2.4 & 68.4 \\
\texttt{1GRM\_2} & 41836 & 83672 & 1 & 11.3 & 509.3 \\
\texttt{1GRM\_2} & 41836 & 83672 & 4 & 7.6 & 255.3 \\
\texttt{1GRM\_2} & 41836 & 83672 & 10 & 4.7 & 121.6 \\
\texttt{1GRM\_4} & 20918 & 41836 & 1 & 21.9 & 776.0 \\
\texttt{1GRM\_4} & 20918 & 41836 & 4 & 14.8 & 431.4 \\
\texttt{1GRM\_4} & 20918 & 41836 & 10 & 9.0 & 228.6 \\
\texttt{1GRM\_8} & 10459 & 20918 & 1 & 41.6 & 1000.0 \\
\texttt{1GRM\_8} & 10459 & 20918 & 4 & 29.0 & 715.7 \\
\texttt{1GRM\_8} & 10459 & 20918 & 10 & 16.5 & 358.1 \\
\texttt{1GRM\_16} & 5229 & 10458 & 1 & 82.9 & 1000.0 \\
\texttt{1GRM\_16} & 5229 & 10458 & 4 & 52.6 & 978.6 \\
\texttt{1GRM\_16} & 5229 & 10458 & 10 & 32.7 & 516.4 \\
\texttt{2OAR\_1} & 519221 & 1038654 & 1 & 0.8 & 74.8 \\
\texttt{2OAR\_1} & 519221 & 1038654 & 4 & 0.6 & 35.2 \\
\texttt{2OAR\_1} & 519221 & 1038654 & 10 & 0.4 & 15.0 \\
\texttt{2OAR\_2} & 259557 & 519326 & 1 & 1.7 & 125.5 \\
\texttt{2OAR\_2} & 259557 & 519326 & 4 & 1.2 & 63.7 \\
\texttt{2OAR\_2} & 259557 & 519326 & 10 & 0.7 & 28.1 \\
\texttt{2OAR\_4} & 129725 & 259662 & 1 & 3.5 & 209.7 \\
\texttt{2OAR\_4} & 129725 & 259662 & 4 & 2.4 & 120.2 \\
\texttt{2OAR\_4} & 129725 & 259662 & 10 & 1.4 & 54.7 \\
\texttt{2OAR\_8} & 64810 & 129830 & 1 & 6.9 & 362.3 \\
\texttt{2OAR\_8} & 64810 & 129830 & 4 & 4.8 & 205.7 \\
\texttt{2OAR\_8} & 64810 & 129830 & 10 & 2.9 & 97.7 \\
\texttt{2OAR\_16} & 32361 & 64914 & 1 & 13.8 & 600.0 \\
\texttt{2OAR\_16} & 32361 & 64914 & 4 & 9.5 & 364.3 \\
\texttt{2OAR\_16} & 32361 & 64914 & 10 & 5.7 & 161.2 \\
\texttt{3SY7\_1} & 221469 & 442946 & 1 & 2.0 & 141.4 \\
\texttt{3SY7\_1} & 221469 & 442946 & 4 & 1.4 & 75.0 \\
\texttt{3SY7\_1} & 221469 & 442946 & 10 & 0.9 & 32.9 \\
\texttt{3SY7\_2} & 110733 & 221472 & 1 & 4.1 & 253.9 \\
\texttt{3SY7\_2} & 110733 & 221472 & 4 & 2.9 & 131.9 \\
\texttt{3SY7\_2} & 110733 & 221472 & 10 & 1.8 & 60.7 \\
\texttt{3SY7\_4} & 55365 & 110736 & 1 & 6.6 & 431.4 \\
\texttt{3SY7\_4} & 55365 & 110736 & 4 & 5.7 & 236.5 \\
\texttt{3SY7\_4} & 55365 & 110736 & 10 & 3.5 & 107.0 \\
\texttt{3SY7\_8} & 27680 & 55368 & 1 & 16.3 & 688.6 \\
\texttt{3SY7\_8} & 27680 & 55368 & 4 & 11.0 & 402.3 \\
\texttt{3SY7\_8} & 27680 & 55368 & 10 & 6.8 & 179.5 \\
\texttt{3SY7\_16} & 13840 & 27684 & 1 & 31.7 & 958.6 \\
\texttt{3SY7\_16} & 13840 & 27684 & 4 & 21.5 & 610.0 \\
\texttt{3SY7\_16} & 13840 & 27684 & 10 & 13.1 & 298.1 \\
\bottomrule
\end{tabular}
\caption{Measured FPS across mesh resolutions and active cutting planes, CPU vs.\ GPU wedge slicing.}
\end{table}
 directly.
\begin{table}[H]
\centering
\footnotesize
\renewcommand{\arraystretch}{0.92}
\begin{tabular}{@{}lrrrrr@{}}
\toprule
Mesh & Vertices & Triangles & Planes ($M$) & CPU FPS & GPU FPS \\
\midrule
\texttt{1GRM\_1} & 83673 & 167346 & 1 & 5.5 & 283.9 \\
\texttt{1GRM\_1} & 83673 & 167346 & 4 & 3.8 & 152.4 \\
\texttt{1GRM\_1} & 83673 & 167346 & 10 & 2.4 & 68.4 \\
\texttt{1GRM\_2} & 41836 & 83672 & 1 & 11.3 & 509.3 \\
\texttt{1GRM\_2} & 41836 & 83672 & 4 & 7.6 & 255.3 \\
\texttt{1GRM\_2} & 41836 & 83672 & 10 & 4.7 & 121.6 \\
\texttt{1GRM\_4} & 20918 & 41836 & 1 & 21.9 & 776.0 \\
\texttt{1GRM\_4} & 20918 & 41836 & 4 & 14.8 & 431.4 \\
\texttt{1GRM\_4} & 20918 & 41836 & 10 & 9.0 & 228.6 \\
\texttt{1GRM\_8} & 10459 & 20918 & 1 & 41.6 & 1000.0 \\
\texttt{1GRM\_8} & 10459 & 20918 & 4 & 29.0 & 715.7 \\
\texttt{1GRM\_8} & 10459 & 20918 & 10 & 16.5 & 358.1 \\
\texttt{1GRM\_16} & 5229 & 10458 & 1 & 82.9 & 1000.0 \\
\texttt{1GRM\_16} & 5229 & 10458 & 4 & 52.6 & 978.6 \\
\texttt{1GRM\_16} & 5229 & 10458 & 10 & 32.7 & 516.4 \\
\texttt{2OAR\_1} & 519221 & 1038654 & 1 & 0.8 & 74.8 \\
\texttt{2OAR\_1} & 519221 & 1038654 & 4 & 0.6 & 35.2 \\
\texttt{2OAR\_1} & 519221 & 1038654 & 10 & 0.4 & 15.0 \\
\texttt{2OAR\_2} & 259557 & 519326 & 1 & 1.7 & 125.5 \\
\texttt{2OAR\_2} & 259557 & 519326 & 4 & 1.2 & 63.7 \\
\texttt{2OAR\_2} & 259557 & 519326 & 10 & 0.7 & 28.1 \\
\texttt{2OAR\_4} & 129725 & 259662 & 1 & 3.5 & 209.7 \\
\texttt{2OAR\_4} & 129725 & 259662 & 4 & 2.4 & 120.2 \\
\texttt{2OAR\_4} & 129725 & 259662 & 10 & 1.4 & 54.7 \\
\texttt{2OAR\_8} & 64810 & 129830 & 1 & 6.9 & 362.3 \\
\texttt{2OAR\_8} & 64810 & 129830 & 4 & 4.8 & 205.7 \\
\texttt{2OAR\_8} & 64810 & 129830 & 10 & 2.9 & 97.7 \\
\texttt{2OAR\_16} & 32361 & 64914 & 1 & 13.8 & 600.0 \\
\texttt{2OAR\_16} & 32361 & 64914 & 4 & 9.5 & 364.3 \\
\texttt{2OAR\_16} & 32361 & 64914 & 10 & 5.7 & 161.2 \\
\texttt{3SY7\_1} & 221469 & 442946 & 1 & 2.0 & 141.4 \\
\texttt{3SY7\_1} & 221469 & 442946 & 4 & 1.4 & 75.0 \\
\texttt{3SY7\_1} & 221469 & 442946 & 10 & 0.9 & 32.9 \\
\texttt{3SY7\_2} & 110733 & 221472 & 1 & 4.1 & 253.9 \\
\texttt{3SY7\_2} & 110733 & 221472 & 4 & 2.9 & 131.9 \\
\texttt{3SY7\_2} & 110733 & 221472 & 10 & 1.8 & 60.7 \\
\texttt{3SY7\_4} & 55365 & 110736 & 1 & 6.6 & 431.4 \\
\texttt{3SY7\_4} & 55365 & 110736 & 4 & 5.7 & 236.5 \\
\texttt{3SY7\_4} & 55365 & 110736 & 10 & 3.5 & 107.0 \\
\texttt{3SY7\_8} & 27680 & 55368 & 1 & 16.3 & 688.6 \\
\texttt{3SY7\_8} & 27680 & 55368 & 4 & 11.0 & 402.3 \\
\texttt{3SY7\_8} & 27680 & 55368 & 10 & 6.8 & 179.5 \\
\texttt{3SY7\_16} & 13840 & 27684 & 1 & 31.7 & 958.6 \\
\texttt{3SY7\_16} & 13840 & 27684 & 4 & 21.5 & 610.0 \\
\texttt{3SY7\_16} & 13840 & 27684 & 10 & 13.1 & 298.1 \\
\bottomrule
\end{tabular}
\caption{Measured FPS across mesh resolutions and active cutting planes, CPU vs.\ GPU wedge slicing.}
\end{table}

\vspace{-4pt}


\section*{Extension: Non Concurrent Cutting Planes}
Tasks 3 and 4 cut with planes that all pass through the scene centre. That
restriction is what keeps the decomposition simple: $M$ concurrent planes always
give exactly $2M$ wedges, each bounded by two half-planes, so the regions can be
worked out in advance (\texttt{rebuildWedgeGeometryDefs}) and each wedge costs the
same two clips no matter how large $M$ is.

Mode~4 (\texttt{4}) removes that restriction. With planes in general position, the
number of regions depends on where the planes happen to fall, so it cannot be
computed up front. \texttt{buildArrangementCPU} therefore finds the regions by
cutting repeatedly: start with the whole mesh, and for each plane split every
surviving piece in two, keeping both halves and discarding anything empty.

Each split also pushes its two halves apart along the normal of the plane that
separated them, and these displacements add up as a piece passes through
successive planes. This is what guarantees separation: two pieces divided by a
plane always move in opposite directions along it, whatever their size. Pushing
each piece along its own centroid direction does not work, because a thin sliver
and the slab next to it have almost the same centroid and end up on top of each
other.

None of this ports to the geometry shader. A GS handles one triangle at a time
and carries no state between passes, so the accumulated offsets have nowhere to
live, and $M$ chained clips would need $3\cdot2^M$ output vertices against a
fixed \texttt{max\_vertices} budget.


\section*{Keyboard Controls}
\begin{tabular}{@{}ll@{}}
\toprule
Key & Action \\
\midrule
\texttt{1}/\texttt{2}/\texttt{3}/\texttt{4} & Whole mesh / CPU slice / GPU slice / non-concurrent planes \\
\texttt{j}/\texttt{k} & Decrease / increase active cutting planes ($M$) \\
\texttt{o} & Toggle cutting planes through-centre / off-centre \\
\texttt{n} & Toggle vertex-normal weighting (area / angle) \\
\texttt{m} & Toggle flat (face-normal) / smooth (vertex-normal) shading \\
\texttt{v} & Cycle colormap (viridis / coolwarm) \\
\texttt{c} & Toggle lit colormap vs.\ flat colour \\
\texttt{f} & Toggle wireframe / filled \\
\texttt{d} & Toggle dark / light background \\
\texttt{+}/\texttt{-} & Zoom in / out \\
\texttt{b} & Run the benchmark sweep \\
\bottomrule
\end{tabular}
\end{document}